\documentclass{article}

\usepackage{amsmath}
\usepackage{enumitem}
\usepackage[boxed,vlined,boxruled]{algorithm2e}

\title{Suffix Array}
\author{Kanishk Goel}
\date{7 May, 2026}
\parindent 0px

\SetAlgoCaptionLayout{centerline}

\begin{document}
	\SetAlgoCaptionSeparator{}
	\renewcommand{\thealgocf}{} 
	\renewcommand{\AlCapNameFnt}{\bfseries}
	\renewcommand{\AlCapFnt}{}
	\DontPrintSemicolon
	\renewcommand{\algorithmcfname}{}
	\maketitle
	\SetAlgoNoEnd
	
	\section{Suffixes}
	For a string abcde, suffixes are abcde, bcde, cde, de and e.
	
	\section{Suffix Array}
	A suffix array contains empty string as well. It is the lexicographical order of all suffixes.
	Storing it would require $O(n^2)$ space, so we store the indexes of the strings at the place which they start in the original string. \newline
	
	For example - ababba
	\begin{enumerate}
		\setcounter{enumi}{-1}
		\item (6)
		\item a (5)
		\item ababba (0)
		\item abba (2)
		\item ba (4)
		\item babba (1)
		\item bba (3)
	\end{enumerate}
	where $(\cdot)$ indicates the numbering of the index i.e the way it is stored in the suffix array. \newline
	
	The algorithm is made simpler by adding a new symbol which is treated as the smallest symbol in place of the empty string, then everything is made of in powers of 2 by adding some filler strings in the suffixes. Each filler character is simply its $index \; mod \; (length \, of \, string)$.
	
	\section{Algorithm Idea}
	Start with base $k = 0$, this means we are comparing all strings by the first character. Use any $O(n \, log \, n)$ sort for this step. The algorithm is divided into $log \, n$ iterations. For $k^{th}$ iteration, the string $s$ contains $s[i \ldots i+ 2^k - 1]$ slice. (Here $i$ is the index of the suffix and is computed $i \, mod \, n$) \newline
	
	To compare string of length $2^{k+1}$, we first compare the number we gave to left part of the string, if lefts are unequal, we don't need to compare the right, else compare right.  \newline
	
	Done quickly through use of \textbf{Prefix Doubling}.
	
	\section{Prefix Doubling}
	
	DP technique to compare $2^n$ length strings. 
	\subsection{Base Case$(2^0)$}
	You compare the length $1$ strings directly and assign equivalence classes. For example a,b,a get 1,2,1.
	
	\subsection{Transition $2^k$ to $2^{k+1}$}
	Divide into $2^{k+1}$ strings and for $k$ the classes have been computed. Compare left half, if equal compare right half. So, the comparision step becomes $O(1)$.
	
	\subsection{The DP}
	For example, banana is the string, for k = 0, we insert into DP[k][i] (null strings are represented as \$) as follows: \newline 
	
	Dp[k=0] = [b,a,n,a,n,a,\$] = [2,1,3,1,3,1,0] \newline 
	
	Dp[k=1] = [ba,an,na,an,na,a\$,\$b] = [3,2,4,2,4,1,0] 
	
	according to the rule of pairs  (DP[$k=1$][$i$],DP[$k=1$][$i+2^{k-1}$]), 
	
	therefore, what was sorted - [(2,1),(1,3) and so on]
	
	
\end{document}